% Canonical Paper IV section.
Ordered binary decision diagrams give a finite model in which contextual
residual classes are represented directly as graph nodes.  The example also
shows that exponential cut sensitivity does not require noncommuting semantic
operations.

\subsection{Residual functions after a variable prefix}

Let $F:\{0,1\}^N\to\{0,1\}$ and fix a variable order $\pi$.  After assigning
the first $t$ variables according to a prefix $\alpha$, the remaining behavior
is the restricted function $F|_\alpha$.

\begin{proposition}[Residual-width lower bound]
\label{prop:obdd-residual-lower-bound}
Every deterministic OBDD respecting $\pi$ has, at the cut after the first $t$
variables, at least as many distinct reachable nodes as there are distinct
restricted functions $F|_\alpha$.
\end{proposition}

\begin{proof}
If two prefixes reach the same node, every common continuation follows the same
remaining subgraph and therefore gives the same Boolean answer.  Hence the two
restricted functions are equal.  Distinct residual functions must reach
distinct nodes.
\end{proof}

For reduced OBDDs the converse also holds: isomorphic remaining subfunctions
are merged.  This is the decision-diagram instance of
Theorem~\ref{thm:minimal-exact-residual}; see~\cite{Bryant1986,Wegener2000}.

\subsection{Equality under block order}

Define
\[
  \operatorname{EQ}_n(x,y)=1
  \quad\Longleftrightarrow\quad
  x=y,
  \qquad x,y\in\{0,1\}^n.
\]
Consider the block order
\[
  x_1,\ldots,x_n,y_1,\ldots,y_n.
\]
After all $x$-variables have been assigned to a string $a$, the residual
function is
\[
  g_a(y)=\mathbf1[y=a].
\]

\begin{theorem}[Exponential middle width]
\label{thm:eq-block-width}
Every OBDD for $\operatorname{EQ}_n$ in block order has width at least $2^n$
at the middle cut.
\end{theorem}

\begin{proof}
The functions $g_a$, $a\in\{0,1\}^n$, are pairwise distinct.  Apply
Proposition~\ref{prop:obdd-residual-lower-bound}.
\end{proof}

The logarithm of this width is $n$ bits, exactly the information needed to
remember the first block for arbitrary comparison with the second.

\subsection{Equality under interleaved order}

Now use
\[
  x_1,y_1,x_2,y_2,\ldots,x_n,y_n.
\]
While equality has not failed, reading $x_i$ creates one of two pending states,
``the next $y_i$ must be $0$'' or ``the next $y_i$ must be $1$''.  Reading
$y_i$ returns to the matching state or goes permanently to the zero terminal.

\begin{theorem}[Constant interleaved width]
\label{thm:eq-interleaved-width}
There is an OBDD for $\operatorname{EQ}_n$ in interleaved order with maximum
width at most three and total size $O(n)$.
\end{theorem}

\begin{proof}
Use the matching state, the two pending-bit states, and the shared zero
terminal as described above.  At any variable level at most three nonterminal
or active residual states are needed, and each pair of variables contributes
only a constant number of nodes.
\end{proof}

The two orders therefore expose exponentially different residual widths for
the same Boolean function and representation language.

\subsection{The restrictions commute}

Assignments to two distinct variables commute as semantic restriction
operators:
\[
  (F|_{x_i=a})|_{x_j=b}
  =(F|_{x_j=b})|_{x_i=a}.
\]
Nevertheless, the order in which variables cross the cut changes which partial
information must remain available.  The exponential separation in
Theorems~\ref{thm:eq-block-width} and
\ref{thm:eq-interleaved-width} is therefore a schedule and representation
phenomenon, not a consequence of operator noncommutativity.

\subsection{What the width theorem measures}

At one cut, $\log_2(\text{width})$ lower-bounds the number of bits needed to
select a residual node.  The entire OBDD also stores transitions for all input
prefixes, so its static node count is a different resource.  Moreover, a
program outside the OBDD model might compute equality with a different state
organization.  The theorem is exact for fixed variable order and the OBDD
representation class; it is not a model-independent lower bound for every
algorithm computing equality.
